% Transcription — fonds Alexandre Grothendieck, shelfmark 45, batch 2 (pages 21-40).
% Established from the facsimile archives/batches/45/batch-02.pdf, cut from a
% copy of 45.pdf fetched through the project's relay
% (https://grothendieck-relay.onrender.com/source/45.pdf); PDF page k =
% archivists' page 20+k, no cover sheet. Per the skill transcribe-grothendieck.
% The folder is seventy-nine pages in four batches (1-20, 21-40, 41-60,
% 61-79); this is the second, the others are done by separate passes.
%
% Pass: Opus 5.5 (claude-opus-5-5), 2026-09-26 — first pass, unchecked against the pages by a human.
%
% Covers, no \page marker: page 27, a grey sheet blank but for « Groupes
% [diagonalisables struck] radiciels et Hyperalgèbres des groupes
% diagonalisables » top right; page 36, a tan sheet blank but for
% « Cohomologie de l'algèbre de Lie et cohom. des groupes ... » top right.
% Both in his hand; their words become section headings, with a \note each.
%
% Pages skipped, without description: 30, 32, 38, 40. They are typed versos
% (reused paper) carrying nothing in his hand: 30 and 32 are stencilled pages
% 21-22 of a text numbered « n° 246 » on ordered planes and half-lines; 38
% and 40 are a typed text on integral preschemes (domaine de définition of a
% rational function, separation). Neither is annotated. Pages summarised:
% none; page 23, a sheet of sketches written in every direction, is given in
% its legible fragments, with a \note saying so.
%
% The 1961 annotated typescript of the folder title does not appear in pages
% 21-40. The four typed versos above are not it: they carry no annotation,
% they are scrap paper under his own notes (the stencil of page 30 shows
% through page 29, the typescript of page 38 through page 37), and neither
% text is about semisimple groups. So no decision on its authorship is made
% in this batch.
%
% Pages transcribed: 21, 22, 24, 25, 26, 28, 29, 31, 33, 34, 35, 37, 39, and
% 23 in fragments. His pagination: none; archivists' numbers only. Page 22,
% the lower part of 21 (a boxed passage), 24 and 26 are cancelled by long
% strokes, recorded once per page. Page 21 does not continue the Théorème
% that page 20 breaks off on; page 24 takes the same Théorème up again
% (H containing a Cartan subgroup, Demazure), in a new form.
%
% Notation, kept from batch 1: gothic g, h, t as \mathfrak{g}, \mathfrak{h},
% \mathfrak{t}; his round L (Lie algebra of a Borel B) as \mathfrak{L};
% Drap(G), Grass(g), Norm_G. New in this batch: U_A(G) the hyperalgebra
% (\mathcal{U}_A(G)), Aff_A(G) the affine algebra, X_n the divided powers
% basis, gamma^n X divided powers, \check{E} the dual, s = t-1.
%
% What the batch is about. Pages 21-26: the Borel subgroup schemes of a
% reductive group scheme over S (is a smooth subgroup with Borel fibres a
% Borel? reduction to the semisimple adjoint case; H^1(B, g/L) = 0 as the
% infinitesimal question), a cancelled Théorème on the radical being defined
% over k, the Théorème that a smooth connected H containing a Cartan subgroup
% is the connected normaliser of its Lie algebra (after Demazure), a
% Théorème B = Norm_G(B), B -> Norm_G(L) surjective, the equivalent
% conditions i)-iii), and on page 26 a finiteness argument for a group K via
% the kernels of n.id_K, with an application to Pic^tau/Pic^0 and the
% constancy of the torsion of Néron-Severi. Pages 28-35: hyperalgebras —
% of the additive group (divided powers, Delta exp X = exp X ⊗ exp Y), of
% V(E), of the multiplicative group (computations with X_n = binom(X,n),
% s = t-1, the structure constants of X_p X_q), then the duality between
% Aff_A(G) and U_A(G) and Propositions 1 and 2 (descent of elements from
% B to A). Pages 37-39: G simple over an Artin local S, the map
% G -> Aut_S(g) and the conditions H^0, H^1, H^2(g_0, g_0) = 0, and the
% question of a spectral sequence linking H^*(G,-) and H^*(g,-).

\documentclass[11pt,a4paper]{article}
% Le préambule commun à toutes les transcriptions du fonds Grothendieck.
%
% Il tient en une idée : **l'incertitude se compose, elle ne se résout pas.**
% Une transcription qui lisse un mot douteux a détruit la seule chose pour
% laquelle on la faisait — un lecteur doit pouvoir savoir, à chaque endroit,
% ce qui était sur le feuillet et ce qui a été ajouté. Les six macros qui
% suivent sont donc l'appareil critique, et non de la décoration.
%
% Les mêmes commandes sont reconnues par `scripts/render.mjs`, qui produit la
% vue de lecture du site. Une macro ajoutée ici doit l'être aussi là-bas,
% faute de quoi le rendu échouera bruyamment — ce qui est voulu : un rendu
% silencieusement faux est pire qu'un rendu absent.

\ProvidesPackage{grothendieck}[2026/08/08 Transcriptions du fonds Grothendieck]

\usepackage{amsmath,amssymb,amsthm}
\usepackage{mathrsfs}
\usepackage{stmaryrd}
% Les diagrammes commutatifs sont partout dans ce fonds : c'est la langue
% même de la géométrie algébrique de Grothendieck, et une transcription qui
% les rendrait en prose ne transcrirait rien.
\usepackage{tikz-cd}
% Français par défaut — c'est la langue des feuillets — mais l'anglais est
% chargé aussi : la traduction et le résumé sont des documents anglais, et la
% typographie française (espaces avant les deux-points) y serait une faute.
% Ces documents basculent par \selectlanguage{english} en tête de corps.
\usepackage[english,french]{babel}
\usepackage{fontspec}
\usepackage{geometry}
\geometry{a4paper,margin=25mm,marginparwidth=28mm}
\usepackage{marginnote}
\usepackage{xcolor}
% ulem, not soul. `soul`'s \st and \ul are built on a character-by-character
% scan that XeLaTeX's font machinery breaks: the document fails with an
% undefined control sequence rather than degrading. ulem does the same two jobs
% and is Unicode-engine safe. `normalem` stops it hijacking \emph, which the
% transcriptions use for Grothendieck's own emphasis.
\usepackage[normalem]{ulem}
\usepackage{hyperref}
\hypersetup{colorlinks=true,linkcolor=black,urlcolor=black,pdfborder={0 0 0}}

% --- Métadonnées du lot -----------------------------------------------------
% Reprises telles quelles par le script de rendu, qui les affiche en tête de
% la vue de lecture. Elles fixent ce que le fichier prétend couvrir : sans
% elles, une transcription flotte sans point d'attache dans un fonds de seize
% mille feuillets.

\newcommand{\folder}[1]{\gdef\grothFolder{#1}}
\newcommand{\batch}[1]{\gdef\grothBatch{#1}}
\newcommand{\leaves}[2]{\gdef\grothFirst{#1}\gdef\grothLast{#2}}
\newcommand{\foldertitle}[1]{\gdef\grothFoldertitle{#1}}
\gdef\grothFolder{?}\gdef\grothBatch{?}\gdef\grothFirst{?}\gdef\grothLast{?}\gdef\grothFoldertitle{}

% --- Le résumé --------------------------------------------------------------
%
% Il ouvre la lecture modernisée, sous le titre « Résumé », et il est composé
% en petit. Ce n'est pas un ornement : c'est le seul passage du document qui
% ne soit pas des mathématiques, et un lecteur doit pouvoir le reconnaître
% avant de l'avoir lu — puis le sauter s'il connaît déjà le sujet, ou s'y
% arrêter s'il ne le connaît pas. Le filet qui le suit marque la reprise du
% corps.

\newenvironment{resume}
  {\small\setlength{\parskip}{0.45em}}
  {\par\medskip\hrule\medskip}

% --- Mots-clés ---------------------------------------------------------------
%
% En anglais, en clôture du résumé de la lecture modernisée : le vocabulaire
% moderne sous lequel chercher ce que les feuillets construisent. Le manifeste
% du site (`npm run manifest`) les extrait pour étiqueter la cote entière —
% cette ligne est la source unique des tags, il n'y a pas de fichier à côté.
\newcommand{\keywords}[1]{\par\smallskip\noindent
  {\footnotesize\textsc{Keywords} --- \textit{#1}}\par}

% --- L'appareil critique ----------------------------------------------------

% Le numéro de feuillet, en marge. C'est l'ancre qui permet de régler un
% désaccord sur une formule en regardant *un* feuillet plutôt qu'une chemise
% de six cents. Le site s'en sert aussi pour tourner les pages du fac-similé
% quand on fait défiler la transcription.
\newcommand{\leaf}[1]{%
  \marginnote{\footnotesize\color{gray!70}\textsf{#1}}%
  \label{leaf:#1}\ignorespaces}

% Illisible : on ne devine pas. Un mot inventé qui se lit comme les autres est
% le pire résultat possible.
\newcommand{\ill}{\textcolor{red!60!black}{[\,\ldots\,]}}

% Lecture douteuse : le mot est donné, et signalé comme incertain.
\newcommand{\uncertain}[1]{\uline{#1}}

% Ajout de l'éditeur — une lettre manquante, un mot sous-entendu.
\newcommand{\add}[1]{\textcolor{blue!50!black}{[#1]}}

% Biffé par Grothendieck lui-même. Ce qu'il a rayé fait partie du document :
% une rature dit souvent où la pensée a changé de direction.
\newcommand{\struck}[1]{\sout{#1}}

% Note du transcripteur, sur l'état matériel du feuillet.
\newcommand{\note}[1]{\par\smallskip{\small\color{gray!60!black}\textit{[#1]}}\par\smallskip}

% Note marginale de Grothendieck, distincte du corps du texte.
\newcommand{\marginal}[1]{\par\smallskip\noindent\hspace{1em}%
  {\small\color{gray!60!black}#1}\par\smallskip}

% --- Titre ------------------------------------------------------------------

\AtBeginDocument{%
  \begin{center}
    {\large\bfseries Fonds Alexandre Grothendieck --- cote n\textsuperscript{o}~\grothFolder}\\[2pt]
    {\itshape\grothFoldertitle}\\[4pt]
    {\small feuillets \grothFirst--\grothLast\ (lot \grothBatch)}\\[2pt]
    {\footnotesize\color{gray!60!black}Transcription établie d'après le fac-similé numérisé
     par l'Université de Montpellier.\\
     Les lectures incertaines sont soulignées, les illisibles notées [\,\ldots\,],
     les ajouts entre crochets.}
  \end{center}
  \vspace{1em}}

\endinput


\folder{45}
\batch{2}
\pages{21}{40}
\dating{1961-[vers 1970]}
\watermark{Édition de démonstration}
\foldertitle{Groupes semi-simples : notes manuscrites (s.d.), tapuscrit annoté (1961)}

\begin{document}

\page{21}

\note{la page ne poursuit pas la phrase sur laquelle s'arrête la page 20 ;
elle s'ouvre sur un point numéroté 3)}

3) $G$ \uncertain{schéma} en groupes \struck{\ill{}} \uncertain{réductif} sur
$S$, $B$ \uncertain{sous-schéma} en groupes \uncertain{lisse}$/S$
\uncertain{invariant} \uncertain{tel} que \uncertain{ses} fibres
\uncertain{soient} des Borel. Est-ce que $B$ est un Borel ?? / Oui, on se
\uncertain{ramène} au cas \uncertain{semi-simple} \uncertain{adjoint} \ldots

\struck{\uncertain{En effet}, on peut \uncertain{supposer} que \ill{}
\add{\uncertain{déployé}} $B$ contient un tore maximal \uncertain{de} $G$.
Mais $\mathfrak{L} = \mathfrak{t} + \sum_{\alpha \in S} \mathfrak{g}_\alpha$,
\ill{} $P$ \ill{} \ill{} « racines \uncertain{positives} ». Donc il contient un
Borel $B'$ ayant \uncertain{même} $\mathfrak{L}$.}
\note{ce passage est pris dans un cadre et barré de trois longs traits
obliques ; le mot entouré en fin de première ligne, relié à « $G$ », n'est
pas lu. Sur la dernière ligne, un « Borel » biffé précède « $B'$ »}

Laissons tomber l'hypothèse $G$ réductif,
\struck{\ill{} \ill{} \ill{} \ill{} \ill{} \ill{} deux tels $B$, $B'$
\ill{} \ill{} \ill{} \ill{} \ill{} \ill{}} \add{\uncertain{loc.}}
\uncertain{de type fini} \ill{} \ill{} \uncertain{propre}
\struck{\ill{}} \ill{} \uncertain{qu.-cpt}. Cela se \uncertain{ramène} aux
fibres \uncertain{géométriques} : la question
\[
H^1(B, \mathfrak{g}/\mathfrak{L}) = 0 \quad ??
\]
\uncertain{du} point de vue \emph{infinitésimal}.

On peut supposer que $B$ et $B'$ contiennent un même tore $T$, et [en se
\uncertain{ramenant} au \uncertain{voisinage} d'un point $s \in S$] que
$B_s = B'_s$.

\page{22}

\note{la page est barrée de trois longs traits obliques ; le texte est
bordé à gauche d'un trait vertical}

\emph{Théorème} (?) $G$ \uncertain{groupe} \uncertain{algébrique} lisse
\uncertain{connexe}$/k$. Soit $\bar k$ la clôture algébrique,
$\bar G = G_{\bar k}$, $\bar T$ un tore maximal de $\bar G$,
\struck{supposons} $\bar R$ le radical de $\bar G$,
\struck{\uncertain{considérons} \ill{} \ill{} \ill{}}
$\bar{\mathfrak{r}} \subset \mathfrak{g}$ les algèbres de Lie, \ill{},
\uncertain{supposons} que les « racines » \add{\uncertain{nulles}} (de
$\bar T$) \uncertain{intervenant} dans $\mathfrak{g}/\bar{\mathfrak{r}}$
\ill{} \ill{} \ill{} $\bar{\mathfrak{r}}$
[\uncertain{p.\ ex.}\ que $\bar T$ opère trivialement sur $\bar{\mathfrak{r}}$]
\struck{[\uncertain{p.\ ex.}\ que \ill{}]} Alors $\bar R$ est défini sur $k$.
\note{les lettres surlignées sont lues $\bar G$, $\bar T$, $\bar R$ et
$\bar{\mathfrak{r}}$ ; le $\bar{\mathfrak{r}}$ gothique est une lecture
douteuse}

\page{23}

\note{feuillet de croquis écrit en tous sens (la page se lit tournée d'un
quart de tour), couvert de flèches, de cadres et de traits de renvoi, en
partie barré. On n'en donne que les fragments lisibles, dans l'ordre du
haut vers le bas}

Radical \add{\uncertain{étale}} défini sur $k$, pour \uncertain{tt}
\uncertain{sous-corps} \ill{} $k$

$\Uparrow$

Une variété infinitésimale \ill{} des Borel \emph{est triviale}
(\uncertain{i.e.} \ill{} \ill{} \ill{} $G$)

$\Updownarrow$

Une telle variété, avec un tore maximal \uncertain{fixé}, $G$ \ill{}
\uncertain{constante}.
\note{« est triviale » et « avec un tore maximal » sont soulignés ; une
insertion entourée au-dessus de la première ligne n'est pas lue}

\[
H^1(\mathfrak{R}, \mathfrak{g}/\mathfrak{L})^{B/R} = 0
\]
($R$ radical \uncertain{unipotent})
\[
H^1(B', \mathfrak{g}/\mathfrak{L}) \longrightarrow
H^1(B, \mathfrak{g}/\mathfrak{L}) \longrightarrow
H^1(\,\cdot\,, \mathfrak{g}/\mathfrak{L})
\]
\note{le premier argument du dernier terme n'est pas lu}
\[
(\mathfrak{g}/[\mathfrak{u},\mathfrak{u}])^T = 0, \qquad
\mathfrak{L}_1 = \mathfrak{L}, \qquad B = B', \qquad B_0 = B'_0
\]
\note{ces formules sont éparses sur la feuille, reliées par des flèches ;
une flèche verticale $G \to G'$ en gros traits est dessinée deux fois, avec
un « $H^1(B,\mathfrak{g}/\mathfrak{L})$ » à côté ; en marge un signe dièse suivi de « 2,3,5 »}

\marginal{il suffit que \ill{} les conditions \ill{} $T$ \uncertain{racines}
\ill{} \uncertain{sur} $\mathfrak{g}/\mathfrak{L}$ \ill{} intervenant
\ill{}}
\note{dans un cadre en bas à gauche}

Cela \uncertain{diffère} \ill{} $G.\mathfrak{g}$ \struck{$sl(2).gl(2)$} pour
$G = Sl(2)$ (\uncertain{car.\ quelconque}).

\page{24}

\note{la page est barrée d'une grande croix. Au-dessus du titre, une ligne
entourée : « \uncertain{$G$ connexe} \ill{} » ; elle reprend le Théorème
de la page 20}

\emph{Théorème} (Soit $H$ un sous-groupe de $G$, lisse$/S$ et à fibres
connexes, ayant \uncertain{un} \uncertain{rang} réductif \ill{} au rang
\uncertain{réductif} \uncertain{de} $G$, \struck{Alors} \struck{$H$ est le
\ill{}} $\mathfrak{h} = \mathrm{Lie}(H) \subset \mathfrak{g} =
\mathrm{Lie}(G)$. Alors $H$ est le \emph{normalisateur connexe} de
$\mathfrak{h}$ dans $G$ \struck{(\uncertain{pour} \ill{} \ill{})}
\uncertain{opérant} sur $\mathfrak{g}$ par la représentation adjointe).
\note{le $G$ de « dans $G$ » est repassé}

En effet, \struck{Demazure} \uncertain{Demazure a montré} qu'il suffit de
prouver que \uncertain{les} \uncertain{fibres} \ill{}
\uncertain{géométriques}, le \uncertain{normalisateur} de $\mathfrak{h}$ dans
$\mathfrak{g}$ est égal à $\mathfrak{h}$. On peut \uncertain{supposer} $S$
spectre d'un corps alg.\ clos, et \struck{\ill{}} \uncertain{la}
\uncertain{hypothèse} \uncertain{signifie} alors que $H$ contient un
sous-groupe de Cartan $C$ de $G$, \uncertain{ou} $T$ un tore maximal.

\page{25}

\emph{Théorème}. $G$ schéma en groupes affine \uncertain{connexe} sur $k$,
$B$ Borel, $\mathfrak{L} \subset \mathfrak{g}$ alg.\ de Lie. Alors
\begin{itemize}
  \item[a)] $B = \mathrm{Norm}_G(B)$
  \item[b)] $B \to \mathrm{Norm}_G(\mathfrak{L})$ est \emph{surjectif},
  \uncertain{donc} $\mathrm{Drap}(G) \to \mathrm{Grass}(\mathfrak{g})$
  injectif (\uncertain{pts} \uncertain{à valeurs} \ill{} \ill{})
  \item[c)] Conditions équivalentes
\end{itemize}
\note{un « ? » en marge, relié par un trait à b) ; un signe biffé en marge
devant c)}
\[
\text{i) } B = \mathrm{Norm}_G(\mathfrak{L}) \qquad
\text{ii) } \mathfrak{L} = \mathrm{Norm}_{\mathfrak{g}}(\mathfrak{L})
\]
iii) $\mathrm{Drap}(G) \to \mathrm{Grass}(\mathfrak{g})$
\uncertain{monomorphisme}
\note{dans « $B = \mathrm{Norm}_G(\mathfrak{L})$ » le signe $=$ est
surchargé ; « $\mathrm{Drap}$ » de iii) est souligné}

On peut supposer \uncertain{$k$} alg.\ clos.
\begin{itemize}
  \item[a)] Réduction au cas \uncertain{semi-simple} \ldots
  \item[b)] \struck{\ill{}} Réduction au cas \uncertain{semi-simple}
  \struck{\ill{}} \ill{}.
\end{itemize}
\struck{Soit $U = B_u$}
\note{« semi- » est écrit avec le signe ½}

\page{26}

\note{la page est barrée de quatre longs traits verticaux ; elle commence
au milieu d'un argument}

D'autre part, le degré de $K_s$ ($s \in S$) étant fonction constructible,
on peut \uncertain{supposer} \uncertain{ce} degré \uncertain{fini} et
\uncertain{constant}, sauf à \ill{} \uncertain{un} multiple commun.
Considérons
\[
K \xrightarrow{\;u_n = n.\mathrm{id}_K\;} K .
\]
\struck{C'est \ill{} l'image $N$ \ill{} \ill{} \ill{}}
Mais ${}_nK = \mathrm{Ker}\, u_n$, \uncertain{on} \emph{\uncertain{la}
\uncertain{section} \uncertain{unité} de $K$ \uncertain{est} \uncertain{un}
\uncertain{immersion} \uncertain{ouverte} \uncertain{dans}} ${}_nK$ est un
\uncertain{sous-schéma} \uncertain{ouvert} \uncertain{et} \uncertain{fermé}
de $K$. \uncertain{Comme} il contient les fibres de $K$ \uncertain{par}
\uncertain{construction}, il est \uncertain{égal} à $K$ \ill{}, $K$
\ill{} \ill{} \ill{} \uncertain{schématiquement}.

\emph{NB}. La condition « \uncertain{être} » \ill{} \uncertain{fini}
\uncertain{sur} \ill{} \ill{} \ill{} \uncertain{à} \uncertain{base}
\add{\uncertain{ex.}\ \uncertain{schémas} de \uncertain{Picard}
\ill{}} \ill{} \ill{} \uncertain{sur} $S$. Une \uncertain{application}
\uncertain{intéressante} : supposons que $X/S$ plat propre,
\struck{\ill{}} $\underline{\mathrm{Pic}}^{\tau}_{X/S}$ \struck{\ill{}}
existant, supposons que \add{\uncertain{défini}} en un pt $s$,
$\underline{\mathrm{Pic}}^{0}_{X_s/k(s)}$ \uncertain{soit} \uncertain{lisse},
\uncertain{alors} $\underline{\mathrm{Pic}}^{0}_{X/S}$ existe au
\uncertain{voisinage} de $s$, \uncertain{et}
$\underline{\mathrm{Pic}}^{\tau}_{X/S}/\underline{\mathrm{Pic}}^{0}_{X/S}$
existe, \uncertain{et} \ill{} \ill{} \uncertain{fini} sur $S$ \ill{}
\uncertain{sur} $S$. Donc [si $\underline{\mathrm{Pic}}^{\tau}_{X/S}$ est
\uncertain{propre} sur $S$] on a la \uncertain{constance} de la torsion de
$NS$.
\note{« Une application intéressante » est encadrée ; les indices des
$\underline{\mathrm{Pic}}$ sont lus $X/S$ et $X_s/k(s)$}

\subsection*{Groupes radiciels et Hyperalgèbres des groupes diagonalisables}

\note{inscrit en haut à droite d'un feuillet par ailleurs blanc (page 27),
de sa main ; « radiciels » est écrit au-dessus de « diagonalisables »
biffé. Le feuillet ne reçoit pas de numéro de page}

\page{28}

Hyperalgèbre du groupe additif
\[
\left\lbrace
\begin{array}{l}
\mathcal{U} = A\{X\} \quad \text{puissances divisées} \\
\Delta \gamma^n X = \displaystyle\sum_{p+q=n} \gamma^p X \otimes \gamma^q X
\end{array}
\right.
\]
\note{entre les deux lignes, une première version biffée :
« $\Delta X^{(n)} = \sum X^{\ldots}$ » ; dans la somme, un premier terme
est biffé avant $\gamma^p X \otimes \gamma^q X$}
i.e.
\[
\Delta \exp X = \exp X \otimes \exp Y
\]
i.e.\ $\exp X \in \widehat{\mathcal{U}}$ est un élément multiplicatif.

$\lambda \mapsto \exp \lambda X$ est une représentation de $A$ dans
$G(\widehat{\mathcal{U}})$.
\note{« $\exp Y$ » est lu tel quel ; le $G$ de $G(\widehat{\mathcal{U}})$
surcharge une autre lettre}

Hyperalgèbre d'un \struck{module localement} fibré vectoriel
$V(\underline{E}) = \mathrm{Spec}\, \mathrm{S}_{\mathcal{O}_S}(\underline{E})$
($\underline{E}$ loc.\ libre \ldots)
\[
\left\lbrace
\begin{array}{l}
\mathcal{U}(\underline{E}) = \mathcal{O}_S\{\check{\underline{E}}\} \quad
\text{puissances divisées sur } \check{\underline{E}} \\
\Delta \gamma^n X = \displaystyle\sum_{p+q=n} \gamma^p X \otimes \gamma^q Y
\end{array}
\right.
\]
i.e.\ $X \mapsto \exp X$ \uncertain{de} $\underline{E}$ dans
$\widehat{\mathcal{U}(\underline{E})}$ \struck{\ill{}} \ill{}
$\underline{E}$ dans $G\widehat{\mathcal{U}(\underline{E})}$ [\uncertain{c'est}
\uncertain{un} \uncertain{homom.}\ de faisceaux en groupes].
\note{« puissances divisées sur $\check{\underline{E}}$ » est souligné ; un
mot est noirci après $\mathcal{O}_S\{\check{\underline{E}}\}$}

Hyperalgèbre du groupe multiplicatif
\note{souligné ; le reste de la page est blanc}

\page{29}

\struck{$\Delta X_n, s$}
\[
\langle \Delta X_n, s^p \otimes s^q \rangle = \langle X_n, s^{p+q} \rangle
= \delta_{n,p+q}
\qquad
\langle X_p . X_q, s^n \rangle = \langle X_p \otimes X_q, (\Delta s)^n \rangle
\]
\marginal{$\Delta s = s \otimes 1 + 1 \otimes s + s \otimes s$}
\note{en haut à droite. Le membre de droite de la seconde formule est
entouré ; une flèche en part vers le calcul qui suit}
\[
= \sum_{\alpha+\beta+\gamma = n} \frac{n!}{\alpha!\,\beta!\,\gamma!}\,
s^{\alpha+\gamma} \otimes s^{\beta+\gamma},
\qquad
\left\lbrace
\begin{array}{l}
\alpha + \gamma = p \\
\beta + \gamma = q
\end{array}
\right.
\]
\[
\left\lbrace
\begin{array}{ll}
0 & \text{si } n \notin [\sup(p,q), p+q] \\[4pt]
\dfrac{n!}{(n-p)!\,(n-q)!\,(p+q-n)!} & \text{sinon}
\end{array}
\right.
\]
\[
\alpha + \beta + 2\gamma = p + q, \qquad
\left\lbrace
\begin{array}{l}
\gamma = p + q - n \\
\alpha = n - q \\
\beta = n - p
\end{array}
\right.
\]
\[
\sup(p,q) \leq n \leq p + q
\]
\note{la première ligne de l'accolade surcharge un premier essai ; « sinon »
est une lecture douteuse}

\[
\boxed{
\begin{array}{l}
\Delta X_n = \displaystyle\sum_{p+q=n} X_p \otimes X_q \\[10pt]
X_p X_q = \displaystyle\sum_{\sup(p,q) \leq n \leq p+q}
\frac{n!}{(n-p)!\,(n-q)!\,(p+q-n)!}\, X_n
\end{array}}
\]
\note{encadré de sa main}

\[
\delta X_n = \sum_{\substack{p+q=n \\ p,q>0}} X_p \otimes X_q
= X_n \otimes X_1 + X_{n-1} \otimes X_2 + \cdots + X_2 \otimes X_{n-1}
+ X_1 \otimes X_n
\]
\note{le premier terme de la somme développée est surchargé}

$\mathcal{U}_n$, \quad $\delta X_n$, \quad $X_2$
\[
X_1^2 = 2X_2 + X_1, \qquad
\left\lbrace
\begin{array}{l}
2X_2 = X_1^2 - X_1 \\
\delta X_2 = X_1 \otimes X_1
\end{array}
\right.
\qquad ? \quad 2\delta X_2 = \delta(X_1^2 - X_1)
\]
\struck{$2\delta(X_1$}
\[
2X_1 \otimes X_1 = \delta(X_1^2 - X_1)
\]
\[
\delta(X_1) = X_1 \otimes 1 + 1 \otimes X_1, \qquad
\delta(X_1^2) = X_1^2 \otimes 1 + 1 \otimes X_1^2 + 2 X_1 \otimes X_1
\]
\note{les signes de ces deux dernières formules sont mal formés ; on les
lit $+$}

\struck{$\delta : \mathcal{U} \to \mathcal{U} \otimes \mathcal{U} \to \mathcal{U}$}
\[
A \longleftarrow A \otimes A \longleftarrow A, \qquad
G \xrightarrow{\ \text{diag.}\ } G \times G \xrightarrow{\ \text{mult.}\ } G
\]
\note{les étiquettes « diag. » et « mult. » sont écrites sous les flèches}

\page{31}

\[
X_n(t^m) = \binom{m}{n} t^m \qquad
\langle X_n, t^m \rangle = \binom{m}{n}
\]
\[
X_n X_{n'}(t^m) = \binom{m}{n}\binom{m}{n'}
\]
\[
\frac{m(m-1)\cdots(m-n+1)}{n!}\;\frac{m(m-1)\cdots(m-n'+1)}{n'!}
\qquad \text{\uncertain{nul} si}
\]
\struck{$X_n, s^{i+c}$}
\[
t = 1 + s, \quad s = t - 1
\]
\[
\langle X_n, s^m \rangle = \langle X_n, (t-1)^m \rangle
\]
\[
= \binom{m}{n} - \binom{m}{1}\binom{m-1}{n} + \binom{m}{2}\binom{m-2}{n}
+ \cdots + (-1)^k \binom{m}{k}\binom{m-k}{n} + \cdots
\]
\note{au milieu de la page, isolés : un grand $X$, « $\Delta X_n$ »,
« $\Delta X -$ »}
\[
\langle X_1, (t-1)^m \rangle = t\,m\,(t-1)^{m-1}\big|_{t=1}
= \left\lbrace
\begin{array}{ll}
0 & \text{si } m \neq 1 \\
1 & \text{si } m = 1
\end{array}
\right.
\]
\[
\langle X_2, (t-1)^m \rangle = \Big\langle \frac{X_1^2}{2} - \frac{X_1}{2} \Big\rangle
\]
\struck{$t-1$}
\[
\frac{m(m-1)}{2!} - \frac{m}{\ \ }
\]
\note{le premier $\langle$ de la dernière ligne est biffé et réécrit}
\struck{$X_n, s^m$}
\[
X_2(t-1) = X_2(t) = \binom{1}{2}
\qquad X_n = \binom{X}{n}
\]
\[
X_n(t^m) = \left\lbrace
\begin{array}{ll}
0 & \text{si } m < n \\
t^m & \text{si } m = n \\
\ldots &
\end{array}
\right. \ (m \geq 0)
\qquad
X_n((t-1)^m) = \left\lbrace
\begin{array}{ll}
0 & \text{si } 0 \leq m < n \\
t^n & \text{si } m = n \\
0 & \text{si } \ldots
\end{array}
\right.
\]
\[
(t-1)^r \qquad \langle X_n, s^m \rangle \qquad
\langle X_n, s^m \rangle = \delta_{nm}
\]
\note{« $m \geq 0$ » est écrit sous l'accolade, précédé d'une flèche
biffée ; la condition « $0 \leq m < n$ » est soulignée}

\page{33}

Soit $G$ un \add{schéma en} groupe affine sur $A$, $A$-simple
\begin{align*}
&\mathrm{Aff}_A(G) \quad \text{son algèbre affine} \\
&\mathcal{U}_A(G) \quad \text{son hyperalgèbre}
\end{align*}
\note{« simple » est ici, comme dans le lot précédent, sa lecture du mot
qui qualifie le schéma sur $A$ ; « schéma en » est ajouté au-dessus de la
ligne}

Donc \struck{\ill{}} $\mathrm{Aff}_A(G)$ est une $A$-algèbre commutative
augmentée, munie d'un homom.\ diagonal multiplicatif, \uncertain{associatif},
compatible avec les augmentations.

$\mathcal{U}_A(G)$ est une $A$-algèbre en général non commutative, munie d'un
homom.\ diagonal multiplicatif, \uncertain{associatif}, \uncertain{commutatif}
avec les augmentations.
\note{« homom. » est une lecture douteuse dans les deux phrases}

Il y a un produit scalaire
\[
\langle f, X \rangle = \langle X, f \rangle
\]
entre $\mathrm{Aff}_A(G)$ et $\mathcal{U}_A(G)$, ce dernier s'identifiant
aux formes $A$-linéaires sur $\mathrm{Aff}_A(G)$ qui s'annulent sur une
puissance de l'idéal d'augmentation.

$\mathcal{U}_A(G)$ opère à gauche et à droite sur $\mathrm{Aff}_A(G)$
[qui \uncertain{est} un bimodule sur $\mathcal{U}_A(G)$]
\[
\left\lbrace
\begin{array}{l}
\langle Y, X * f \rangle = \langle Y * X, f \rangle = \langle f, Y * X \rangle
= \langle f * Y, X \rangle \\
\langle \varepsilon, X * f \rangle = \langle \varepsilon, f * X \rangle
= \langle X, f \rangle
\end{array}
\right.
\]

\page{34}

$\mathrm{Aff}_A(G)$ opère \struck{à gauche et à droite} sur
$\mathcal{U}_A(G)$ ; la formule \uncertain{est} \struck{transposée}
analogue
\[
\langle g, f \lrcorner X \rangle = \langle gf, X \rangle
\qquad \text{en particulier } \langle f, X \rangle = \varepsilon(f \lrcorner X)
\]

Supposons maintenant que $A \subset B$, alors
\[
\mathrm{Aff}_B(G) = \mathrm{Aff}_A(G) \otimes_A B \supset \mathrm{Aff}_A(G)
\]
[on a une inclusion, \uncertain{car} $\mathrm{Aff}_A(G)$ plat sur $A$]
\[
\mathcal{U}_B(G) = \mathcal{U}_A(G) \otimes_A B \supset \mathcal{U}_A(G)
\]
[raison analogue ; la première égalité est due au fait que $G$ est
$A$-simple]

\emph{Proposition 1}. Soit $X \in \mathcal{U}_B(G)$. Conditions équivalentes
\begin{itemize}
  \item[i)] $X \in \mathcal{U}_A(G)$
  \item[ii)] $X * \mathrm{Aff}_A(G) \subset \mathrm{Aff}_A(G)$ \qquad
  ii bis) $\mathrm{Aff}_A(G) * X \subset \mathrm{Aff}_A(G)$
  \item[iii)] $\langle X, \mathrm{Aff}_A(G) \rangle \subset A$
\end{itemize}

En effet, i) $\Rightarrow$ ii) $\Rightarrow$ iii) \add{et i) $\Rightarrow$
ii bis) $\Rightarrow$ iii)} triviaux. iii) implique que \struck{$X$
\ill{}} \struck{\ill{} $X$ \ill{}} $X$ \uncertain{est} une forme
$A$-linéaire sur $\mathrm{Aff}_A(G)$ nulle sur une puissance de l'idéal
d'augmentation, donc un élément $X_0$ de $\mathcal{U}_A(G)$, qui induit
\uncertain{évidemment} \uncertain{défini} par $X$.
\note{« i) $\Rightarrow$ ii bis) $\Rightarrow$ iii) » est écrit au-dessus
de la ligne, avec des flèches doubles ; la phrase s'arrête au bas du
feuillet}

\page{35}

\emph{Proposition 2}. Soit $f \in \mathrm{Aff}_B(G)$
\add{\uncertain{Supposons} $G$ $A$-connexe}. Conditions équivalentes
\begin{itemize}
  \item[i)] $f \in \mathrm{Aff}_A(G)$
  \item[ii)] $f \lrcorner\, \mathcal{U}_A(G) \subset \mathcal{U}_A(G)$
  \item[iii)] $\langle f, \mathcal{U}_A(G) \rangle \subset A$
\end{itemize}
\note{« Supposons $G$ $A$-connexe » est ajouté au-dessus de la ligne, et
remplace un mot biffé}

Évidemment i) $\Rightarrow$ ii) $\Rightarrow$ iii) [et iii) $\Rightarrow$
ii), \uncertain{car} \struck{\ill{}} si $X \in \mathcal{U}_A(G)$, pour
prouver $f \lrcorner X \in \mathcal{U}_A(G)$, il suffit \add{(prop.~1)} de
vérifier que $\langle g, f \lrcorner X \rangle \in A$ pour
$g \in \mathrm{Aff}_A(G)$, or
\[
\langle g, f \lrcorner X \rangle = \langle gf, X \rangle
= \langle f, g \lrcorner X \rangle \in \langle f, \mathcal{U}_A(G) \rangle
\subset A \ ].
\]
Prouvons iii) $\Rightarrow$ i), en utilisant que $G$ est $A$-connexe.
Soit $C = \mathrm{Aff}_A(G)$, considérons le diagramme

\begin{tikzcd}
C \arrow[r] \arrow[d] & C_B \arrow[d] \\
\widehat{C} = \mathrm{Hom}_A(\mathcal{U}, A) \arrow[r] &
\widehat{C}_B = \mathrm{Hom}_B(\mathcal{U}_B, B)
\end{tikzcd}

où ${}^\wedge$ désigne le \add{\uncertain{séparé}} complété pour l'idéal
d'augmentation.
\note{le $C$ de « Soit $C$ » surcharge un $B$ noirci}

\subsection*{Cohomologie de l'algèbre de Lie et cohom.\ des groupes \ldots}

\note{inscrit en haut à droite d'un feuillet par ailleurs blanc (page 36),
de sa main ; le feuillet ne reçoit pas de numéro de page}

\page{37}

Soit $G$ un \uncertain{schéma} en groupes \uncertain{simple} sur
$S = \mathrm{Spec}\, A$, $A$ Artinien local.
\[
\mathfrak{g} = L_S(G)
\]

(1) Si $H^0(\mathfrak{g}_0, \mathfrak{g}_0) = 0$, $\Longrightarrow$
$G \to \underline{\mathrm{Aut}}_S(\mathfrak{g})$ est un
\uncertain{monomorphisme} [donc si $G_0$ n'a pas de sous-groupe
\struck{\uncertain{distingué}} \struck{\ill{}} \uncertain{non} trivial,
alors c'est une \uncertain{immersion} \uncertain{fermée}]
\note{le « (1) » n'est pas écrit ; la ligne commence par « Si », et la
suite de la page renvoie à « (1) et (2) »}

(2) \add{« (1) et »} \struck{\ill{}} $H^1(\mathfrak{g}_0, \mathfrak{g}_0) = 0$
$\Longleftrightarrow$ \struck{\ill{}} $G \to
\underline{\mathrm{Aut}}_S(\mathfrak{g})$ est étale, [donc sous
l'hypothèse dite plus haut, c'est une immersion ouverte, donc
\[
G \simeq \underline{\mathrm{Aut}}^0_S(\mathfrak{g})
\]
si $G$ \ill{} \uncertain{à} fibres connexes]
\note{la flèche de (1) est une flèche double surchargée, celle de (2)
une flèche double réécrite ; dans (2), « (1) et » est ajouté au-dessus
d'un mot biffé}

\marginal{(1) \struck{\uncertain{(1) et (2)}} $\Longrightarrow$
$\underline{\mathrm{Aut}}_S(\mathfrak{g}) \simeq \underline{\mathrm{Aut}}_S(G)$
si $G$ \uncertain{connexe} \ldots\ \uncertain{suffirait} \uncertain{si}
Tohoku \ill{}, Tohoku \struck{\ill{}} \uncertain{vrai} \uncertain{sur}
$\mathbb{Z}$ \ldots\ \uncertain{Utilisant} \uncertain{contre-ex.} \ldots}
\note{note écrite en biais dans la marge gauche, en face de (2) ; lecture
incomplète. « Tohoku » est lu deux fois}

Supposons non plus donné un schéma en groupes $G$, mais seulement
\uncertain{une} $\mathfrak{g}$, considérons
$H = \underline{\mathrm{Aut}}_S(\mathfrak{g})$. \struck{Alors} C'est un
schéma en groupes sur $S$, et
$\mathrm{Lie}_k(H_0) = \mathfrak{h}_0$ est isomorphe à \add{l'algèbre de Lie
des} \struck{\ill{}} \struck{\ill{}} dérivations de $\mathfrak{g}_0$, d'où
$\mathfrak{g}_0 \to \mathfrak{h}_0$, qui est un isom.\ si (1) et (2) sont
\uncertain{satisfaites}. S'il en est ainsi, on ne peut pas affirmer
\uncertain{toutefois} que $H$ est simple sur $S$, \ill{} \ill{} \ill{} \ill{}
\uncertain{suivante}, \struck{\uncertain{Mais} \uncertain{dans}} \ill{}
\ill{} \ill{} \ill{}, \uncertain{si} \ldots
\[
(3) \quad H^2(\mathfrak{g}_0, \mathfrak{g}_0) = 0
\]
\uncertain{qui} \uncertain{signifie} \ill{} \uncertain{aucune}
\add{\uncertain{infinitésimale}}, \uncertain{qu'il} \uncertain{n'y} a
\uncertain{pas} de variation de structure \uncertain{non triviale} de
$\mathfrak{g}_0$, et \emph{implique} que \struck{\ill{} \ill{}}
$\underline{\mathrm{Aut}}_S(\mathfrak{g})$ est \uncertain{simple}
\uncertain{sur} $S$.

\page{39}

\struck{De} \uncertain{plus}, (1) (2) (3) [\uncertain{implique} \ill{},
\uncertain{inversement}, \uncertain{si} \uncertain{on} $G$ \struck{\ill{}}
\uncertain{sur} $S$ \uncertain{donnant} $\mathfrak{g}$] \uncertain{provient}
\uncertain{aussi} [\uncertain{par} $G$ \uncertain{connexe}] \uncertain{de}
$G$ \ldots\ \ill{} \uncertain{de} \uncertain{variations}
\uncertain{infinitésimales} \uncertain{non} \uncertain{triviales}
(\uncertain{Gerstenhaber}),

\uncertain{Question} \uncertain{de} \uncertain{relations} \uncertain{de} la
structure $G$, \uncertain{équivalente} à la question \uncertain{de}
\uncertain{relations} de la structure $\mathfrak{g}$. \uncertain{On}
\uncertain{trouve} \uncertain{une} \uncertain{interprétation} des
$H^3(\mathfrak{g}, \mathfrak{g})$. \ill{}
\note{le $H^3$ est lu sur un signe qui ressemble à « HB » ; la lecture de
l'exposant est douteuse, de même dans le point (3) qui suit}

(3) Si $H^3(\mathfrak{g}, \mathfrak{g}) = 0$, \uncertain{alors} \ill{}
\uncertain{peut} \uncertain{relever} \ill{} \uncertain{groupe}
\struck{\ill{}} \uncertain{unique} $\mathfrak{g}$, \ill{} $G$ \ill{} \ill{}
\uncertain{équivalente}.
\note{un « $H^3$ » isolé est écrit dans la marge en face de (3)}

\emph{Question}. Définir
\[
H^*(G, \mathfrak{g}) \longrightarrow H^*(\mathfrak{g}, \mathfrak{g})
\]
[plus \uncertain{généralement}, si \ill{} \ill{} une représentation linéaire
de $G$ \ldots]. Suite spectrale reliant $H^*(G, \ )$ et
$H^*(\mathfrak{g}, \ )$ ?? Cela exprimerait \uncertain{la}
« \uncertain{différence} » \uncertain{entre} le global \uncertain{et}
l'infinitésimal.

Supposons $G$ semi-simple adjoint. A-t-on alors
\[
H^*(G, \mathfrak{g}) = H^*(\mathfrak{g}, \mathfrak{g}) = 0 \quad ??
\]
\note{« Question » est souligné}

\end{document}
